%-------------------------------------------------------------------------------
\chapter{Physics Modelling}
\label{sec:modelingCorr}
%-------------------------------------------------------------------------------

\todo{This whole section comes from W pT and has to be reworked for mW}

The main signal event samples for $W$ and $Z$ production are generated
using the \POWHEG\ event
generator~\cite{Nason:2004rx,Frixione:2007vw,Alioli:2010xd,Alioli:2008gx}
using the CT10~\cite{Lai:2010vv} PDF, interfaced to \PYTHIA[8]~\cite{Sjostrand:2007gs} using the
AZNLO tune~\cite{STDM-2012-23}. These \PowPythia\ samples are interfaced
to \Photos~\cite{Golonka:2005pn} to simulate the effect of final state
QED radiation. Further details on these samples, on the background simulated samples and on alternative simulated signal samples are in Section~\ref{sec:samples}.

The inclusive \Zll and \Wlnu cross sections are generally
six-dimensional due to the two final-state leptons with fixed mass. As
successfully employed in the 7~\TeV\ \mW\
measurement~\cite{STDM-2014-18,Andari:1966965}, it can be conveniently parameterised
as four quantities related to the di-lepton (or boson) system, one of
which is a trivial $\varphi$ angle in the lab system, and two lepton
decay angles. Instead of a five-dimensional map, it has been shown
that the cross section can be conveniently decomposed to control the
physics modelling in the following way:
\begin{equation}
\label{eqn:master_formula}
\frac{d\sigma}{dp_1 dp_2}=\left\lbrack \frac{d\sigma(m)}{dm} \right\rbrack\left\lbrack \frac{d\sigma(y)}{dy} \right\rbrack\left\lbrack \left.\frac{ d\sigma(\pt)}{d\pt} \right|_y \left(\frac{d\sigma(y)}{dy}\right)^{-1}\right\rbrack \left\lbrack (1+\cos^2\theta) + \sum_{i=0}^7 A_i(\pt,y)P_i(\cos\theta,\phi) \right\rbrack.
\end{equation}

Here $p_1$ and $p_2$ are the lepton and anti-lepton four-momenta; m,
\pt, and y are the invariant mass, transverse momentum, and rapidity
of the dilepton system; $\theta$ and $\phi$ are the polar angle and
azimuth of the lepton\footnote{Here, lepton refers to the negatively
  charged lepton from a \Wminus\ or \Zboson\ boson, and the neutrino
  from a \Wplus\ boson} in any given rest frame of the dilepton
system; $A_i$ are numerical coefficients, and $P_i$ are spherical
harmonics of order zero, one and two.

Such a decomposition specifically allows a ``simple'' modelling of the
mass-distribution (if required), the use of fixed-order predictions
for the $y$ distribution and the 8 $A_i(\pt,y)$ coefficients (with the
analytically known $P_i(\cos\theta,\phi)$ from the spin-1 exchange)
and resummed or tuned parton-shower MC predictions for the boson \pT
at fixed rapidity $d\sigma(\pt)/d\pt|_y$.

\section{Boson Mass}

\todo{Needs to be added}

\section{Boson $p_T$ and rapidity}

\todo{This text comes from the W pT result and needs to be modified for the mW analysis}

The initial agreement between the detector-level transverse momentum
distributions in data and MC is found to be not fully satisfactory,
i.e. the difference is found to be significantly larger than the
experimental sensitivity especially at $13\,\tev$. 
To optimise the analysis results, the MC simulation is already
reweighted by default in \pttruthv to guarantee the best possible data
to MC agreement on the reconstructed \ut (or \ptll) distributions for
$W$ (and $Z$) analyses~\footnote{To achieve this, the full analysis
  including all calibrations and background estimates had to be
  performed already once. Object calibrations, including the Hadronic
  Recoil and related reweightings aer performed such that they are to
  first order insensitive to the boson \pt modelling in the MC.}. This
procedure ensures that the unfolding result is unbiased, and is
referred to as the \texttt{\pt correction}. The uncertainties will be
discussed in later sections of this note and will be derived by
performing as many variations as possible on this correction.

For the $W$ analyses, the agreement between data and
simulation is optimised at detector level for the \ut distribution by
minimising the following $\chi^2$:
\begin{eqnarray}
  \chi^2  &=&  \sum_{ij} \, \Delta^T_i \, C^{-1}_{ij} \, \Delta_j \,   \label{eqX2}
 ,\\
  \Delta_i  &=& ( D_i \, - B_i \,) - \, \sum_{j} \, R_{ij} \times (w_{T}(p_{T}^{W}))_j \, .
\end{eqnarray}
Here $C^{-1}_{ij}$ denotes the total covariance matrix, including statistical and systematic uncertainties on the observed and simulated \ut distributions;
 $\Delta$ is the difference between the background-subtracted data distribution, ($D-B$), and the corresponding reconstruction-level simulated distribution obtained from the product between a generator-level weighting function and the response matrix, $R \, (w_{T}(p_{T}^{W}))$ .
The weighting function, $ w_{T}(p_{T}^{W}) $, evaluated in each truth bin $j$, modifies the generated $W$-boson \pt
distribution in the MC sample used to determine the response matrix for events passing both reconstructed-level and
fiducial volume selection cuts (T\&R).
 The parameters of the $w_{T}(p_{T}^{W})$ weighting function are the degrees of freedom in the minimisation procedure. A satisfactory agreement is obtained using
\begin{equation}
(w_{T}(p_{T}^{W}))_j  = N_j \left[\left(1 + a \, p_{T, j}^{W} + b \, {p_{T, j}^{W}}^2
\right)*\left(1-c+c*r_\mathrm{NNPDF/CT10}(p_{T, j}^{W})\right)\right] \,    \label{eq:Wnom_bin}
\end{equation}
where $r_\mathrm{NNPDF/CT10}(p_{T}^{W})$ represents the ratio between predictions obtained with NNPDF3.0 and CT10NLO in
the full phase space illustrated in Figure~\ref{fig:ratioNNPDFCT10}, and the additional terms provide more flexibility
to the fit. $N_{j}$ is the number of events in truth bin $j$. CT10NLO is very close to the CT10 PDF used in the \POWHEG\ sample, and NNPDF3.0 is representative of more modern PDF sets and was specifically used in the \Sherpa sample. 

The reweighting functions are optimised separately for each process and each centre-of-mass energy. They are derived over the range $\pt = 0 - 100 \gev$ and the value of the correction is frozen at $\ptw=100$~GeV for any larger \pt. The 13~TeV $W$ reweighting functions are shown in \todo{Figure X} for the electron and muon channels. The corresponding plots at $\sqrt{s} = 5$~\TeV\ are shown in \todo{Figure X} . Since the $e$ and $\mu$ channel compatibility is good, the average of the two reweighting functions is eventually used.
The effect of this \pt\ reweighting on the reconstructed \ut distributions is shown in \todo{Figure X}  for 13 TeV and Figure \todo{Figure X}  for 5 TeV.


For the \Zboson\ process the same approach is used, with the addition that separate reweightings are derived for the \ptdilep\ and \ut\ distributions.  Figure~\ref{fig:rew_sf_zpt} shows the resulting best reweighting functions at $\sqrt{s} = 13$~\TeV\ and $\sqrt{s} = 5$~\TeV. For the \ut measurement, the reweighting function is obtained from a simultaneous fit to the electron and muon channels, while  the \ptdilep measurement keeps the averaged reweighting function, as done for the $W$-boson channels (more details are given in Appendix ~\ref{sec:zpt_bias_new}). The general consistency between the reweighting functions obtained for the \ut and \ptdilep distributions, at both center-of-mass energies, indicates that this method catches genuine physics effects. At the same time, it's clear that especially at $\sqrt{s} = 13$~\TeV\ the \ptdilep\ reweighting is able to catch a feature at smallest \pT\ that is lost in the \ut\ resolution.~\footnote{The potential uncertainties of such effects are covered by alternative reweighting functions as is explained later.} Figures~\ref{fig:datamc_nominalMC5} ands~\ref{fig:datamc_nominalMC13} show the reconstruction-level data/MC ratios obtained from this procedure.


The corresponding uncertainties are discussed in \todo{sections X}.


\section{Angular correlations}

The polarisation of the vector boson is entirely encoded in the 8
$A_i$ coefficients of equation~\ref{eqn:master_formula}.  The angular
correlations between the two decay leptons, arising from the
polarisation of the considered vector boson, are not properly taken
into account in our baseline simulation of \Wboson\ and \Zboson\
samples that uses \PowPythia.  This generator is NLO accurate in QCD,
and also suffers from a known feature in the limit where \pttruthv\
goes to 0 ($A_0$ becomes negative). It was shown in
Ref.~\cite{STDM-2014-10} that NNLO predictions for the $A_i$
coefficients correctly describe the measurements for the \Zboson\
at 8~\TeV. Therefore, as in the \mW\ analysis at
7~\TeV~\cite{STDM-2014-18}, the angular coefficients are reweighted
to NNLO predictions.

To obtain this better description of the angular coefficients, DYTURBO is used at NNLO to predict these as a function of boson \pt\ and rapidity for all 3 processes (\Wplus, \Wminus\ and \Zboson). The CT10NNLO PDF set is used.

The $\cos(\theta)$ and $\phi$ are then reweighted according to the procedure described in~\cite{STDM-2014-18}, with the following event weight :
\begin{equation}
\label{eqn:weight_pol}
w = \frac{1+\cos^2\theta + \sum_{i=0}^7 A_i'(\pt,y)P_i(\cos\theta,\phi)}{1+\cos^2\theta + \sum_{i=0}^7 A_i(\pt,y)P_i(\cos\theta,\phi)},
\end{equation}
where angular coefficients of \PowPythia\ are used in the denominator and those predicted at $\mathcal{O}(\alpha_S^2)$ by DYTURBO in the numerator.

The angular coefficients in DYTURBO at NLO, NNLO and in \PowPythia\ are shown as a function of \pttruthv\ in figures \todo{Add any figures}.
